% a01_H0rel part 6x1: the four replaced steps of Balaban's proofs at block factor 5 (input from a01_H0rel_p6.tex)
\paragraph{The four replaced steps at block factor $5$.}
Scale-$i$ units (the block lattice of scale $i$ has spacing $1$); $u_i\defeq1+\ln\ell_i^{-1}$, so $u_{i-1}=u_i+\ln L$, $r_i\defeq r(\ell_i)=Ru_i^{\varrho}$, $p_i=p(\ell_i)=b_0u_i^{p}$ with $\varrho=2$, $p=6$; $s_b=LM$ is the large-block side; $\rho_i\defeq(r(\ell_{i-1})-s_b)/L$ is one previous-scale corridor in scale-$i$ units (\cite[p.~578]{Ba2}); $X'$ denotes the blocks of $X$ read as points of the next lattice \cite[(1.16)]{Ba1}; $\Omega\defeq B^i(\Lambda_2^{(i-1)\prime})$ is the propagator domain and $\Omega_1\defeq B^i(\Lambda_7^{(i-1)\prime})$ the vertex set of \cite[(1.1)--(1.5)]{Ba3}; $\delta_0$, $c_{\mathrm{Ba}}$ and the depth $R_{\mathrm{Ba}}$ are the constants of \cite[Props.~2.1--2.3]{Ba1}. Everything below is a statement about Ba{\l}aban's regions, cut-offs, boxes and partitions, which are built from $(L,M,R,\ell_i,\bar n)$ and the set of event blocks alone; nothing depends on the matter boundary condition, so each statement holds on the cube exactly as on the torus.

\begin{lemma}[the supply of depth, and the criterion]\label{lem:a01:X0}
(a) $\dist(\Lambda^{(i-1)}_{m+s},(\Lambda^{(i-1)}_m)^c)>s\,r(\ell_{i-1})-(s-1)$ for $s\ge1$ (scale-$(i-1)$ units), and in scale-$i$ units $\dist(B^i(\Lambda^{(i-1)\prime}_{m+n}),B^i(\Lambda^{(i-1)\prime}_m)^c)\ge n\,r(\ell_{i-1})/L-2$ for $n\le L+1$. The hypothesis of \cite[Prop.~1]{Ba3} asks depth $>2r_i$ of $\supp\tilde A$ in $\Omega$; at its use in \cite{Ba2}, where the domain is $B^i(\Lambda_0^{(i)})$, the regions supply $\dist(\supp\tilde A,\Omega^c)>2r_i-2-M_\theta-2\eta$ (the slice of $\theta_i$ and the fine blocks), which is all that the proof consumes. (b) Let $\alpha,\beta>0$, $c\ge0$. If $\alpha L\le\beta$, then $\beta\rho_i-\alpha r_i\ge\alpha R\varrho\ln L\,u_i^{\varrho-1}-\beta s_b/L$ for every $i$, so $\beta\rho_i\ge\alpha r_i+c$ at every scale once $R\ge(c+\beta s_b/L)/(\alpha\varrho\ln L)$. If $\alpha L>\beta$, then $\beta\rho_i<\alpha r_i$ for every $\ell_i<\exp\big(1-\ln L/((\alpha L/\beta)^{1/\varrho}-1)\big)$, whatever $R$ and $c$ are. If an object of size $\theta r_{i+1}$ at scale $i+1$ must sit at depth $\beta r_i$, the containment holds at every scale iff $L\theta\le\beta$, and for $L\theta>\beta$ it fails for every $\ell_i<\exp\big(1-\ln L/(1-(\beta/(L\theta))^{1/\varrho})\big)$ (now $r_{i+1}/r_i=(1-\ln L/u_i)^\varrho$ enters). (c) Where only exponential smallness across scales is used, $e^{-\delta\beta\rho_i}\le Ce^{-(\delta\beta/L)r_i}$, which is $O(\ell_i^{\kappa})$ for every $\kappa$ since $\varrho>1$.
\end{lemma}
\begin{proof}
(a) By \cite[(2.8)]{Ba2} $(\Lambda_{m+1})^c$ contains every large block within $r$ of $(\Lambda_m)^c$, so a point of $\Lambda_{m+1}$ is at distance $>r$ from $(\Lambda_m)^c$; along a lattice geodesic from $x\in\Lambda_{m+s}$ to $(\Lambda_m)^c$ each further layer adds $>r-1$. For the primed sets, points of $B^i(y)$, $B^i(y')$ with $y\in\Lambda_{m+n}$, $y'\in\Lambda_m^c$ differ coordinatewise by $(y-y')+\tau$, $|\tau|<L$, so $|x-x'|_\infty>n\,r(\ell_{i-1})-(n-1)-L$, and dividing by $L$ gives the claim for $n\le L+1$. (b) If $\beta/L\ge\alpha$ then $\beta\rho_i-\alpha r_i\ge\alpha(r_{i-1}-r_i)-\beta s_b/L$ and $r_{i-1}-r_i=R[(u_i+\ln L)^\varrho-u_i^\varrho]\ge R\varrho\ln L\,u_i^{\varrho-1}$. If $\alpha L>\beta$, $\beta\rho_i\le(\beta/L)r_{i-1}$ and $r_{i-1}/r_i=(1+\ln L/u_i)^\varrho<\alpha L/\beta$ exactly when $u_i>\ln L/((\alpha L/\beta)^{1/\varrho}-1)$. (c) $e^{-\delta\beta\rho_i}\le e^{\delta\beta s_b/L}e^{-(\delta\beta/L)r_{i-1}}$ and $e^{-cR(1+\ln\ell^{-1})^\varrho}\le\ell^\kappa$ for small $\ell$ because $\varrho>1$.
\end{proof}

Every place where $L$ enters \cite{Ba1,Ba2,Ba3} and \cite[Thm.~3.1]{K1} is a constant, a rate, a convergent sum, a condition on $(b_0,R,M,\eps_0^{\mathrm{Ba}})$, a divisibility, or a \emph{cross-scale containment} $\beta\rho_i\ge\alpha r_i+c$, and by \cref{lem:a01:X0}(b) a containment holds at every scale iff $\alpha L\le\beta$. At $L=5$ exactly four printed steps fail this test (\cref{thm:a06:BALA}): the cube step of \cite[p.~433]{Ba3} ($\alpha=\frac32$, $\beta=5$), the regularity of the vector background \cite[Lemma~2.3]{Ba2} ($\alpha=1$, $\beta=3$), the boxes of \cite[Lemma~2.4]{Ba2} ($\alpha=4$, $\beta=5$; this one fails at every $L\ge2$), and the split \cite[(2.89)]{Ba2} ($L\theta=L$, $\beta=2$); the partition \cite[(3.43)--(3.45)]{Ba2} is a divisibility entry which fails on some lattices at $L=5$. These are gaps of the printed proofs, not counterexamples to the theorems, whose statements never mention the objects concerned.

\subparagraph{(a) The cube step.} On \cite[p.~433]{Ba3} the localisation cubes of a renormalised class, which lie in a cube of side $\ell_{\bar n}$ depending on $\bar n$ only \cite[p.~432]{Ba3}, are enclosed in concentric cubes $\square_1\subset\square$ of sides $r_i$, $3r_i$, sums of large blocks; $G_i(\Omega,\tilde B)$ is written $G_i(\square,\tilde B)+\delta G_i(\Omega,\square,\tilde B)$ with the difference treated as an external field, the regular field $\tilde B$ is replaced on $\square$ by its constant part $B_0$, $B_0$ is gauged away, and $G_i(\square,0)$ is replaced by $G_i(0)$. The step presupposes $\square\subset\Omega$ (resp.\ $\square\subset\Omega_2$): the Neumann operator $G_i(\square,\tilde B)$ is built from the bonds of $\square$, where $\tilde B$ lives only if $\square\subset\Omega$, and \cite[(2.26)]{Ba1} is stated for nested domains. Every vertex of every class lies within sup-distance $1$ of $\Omega_1\cup\supp\tilde A$ (the vertices \cite[(1.6)--(1.7)]{Ba3} are summed over $\Omega_1$; those of \cite[(1.8)--(1.11)]{Ba3} carry $g_i$, $\supp g_i\subset\Omega_1$, or $\tilde A$; the pair vertices \cite[(1.12)--(1.15)]{Ba3} carry $A'$ or $\tilde A$ along a contour inside one block). For an old vertex $x_v$ of the class the containment reads
\begin{equation}\label{eq:a01:N}
\dist(x_v,\Omega^c)>\tfrac32r_\square+2s_b+\ell_{\bar n}\ \Longrightarrow\ \square\subset\Omega\ \text{and}\ \dist(\square_1,\square^c)\ge r_\square
\end{equation}
(cubes rounded outward to large blocks; triangle inequality). With the printed $r_\square=r_i$ and a vertex anchored at $\Omega_1$, whose depth is $\ge5r(\ell_{i-1})/L-2$ by \cref{lem:a01:X0}(a), this is a containment with $\alpha=\frac32$, $\beta=5$: it holds at every scale iff $L\le3$, and at $L=5$ it fails for all small $\ell_i$; the failing vertices ($\lambda$- and $\dm$-vertices summed over $\Omega_1$, $g_i$-localised photon legs near $\partial\Omega_1$) carry no $\tilde A$-leg, so no hypothesis on $\supp\tilde A$ helps.

\begin{proposition}[the cube step with a constant cube]\label{prop:a01:P12}
Let $\ell^\sharp\defeq3\ell_{\bar n}+4N_B$, where $N_B=n_1+n_2-1$ is the number of insertions per line in the expansion of \cite[p.~433]{Ba3} to order $\bar n$ (a function of $\bar n$ only), let $r^\flat$ be $\max\{3\ell^\sharp+R_{\mathrm{Ba}}+12,\,3R_{\mathrm{Ba}}+15,\,107\}$ rounded up to a multiple of $s_b$, and $D^\flat\defeq\frac32r^\flat+2s_b+\ell^\sharp$. Run the step with concentric cubes $\square_1,\square$ of sides $r^\flat$, $3r^\flat$ centred at the component cube. If every old vertex of the class has depth $>D^\flat$ in $\Omega$ (resp.\ $\Omega_2$), the proof of \cite[pp.~433--445]{Ba3} goes through and gives \cite[(1.33)]{Ba3} with a constant of the printed form, depending on $(\alpha_0,\bar n,d,a,M)$ only. If $R\ge R^\flat\defeq D^\flat+3$ and $R\ge M_\theta+2$, then at $L=5$ every class occurring in the use of \cite[Prop.~1]{Ba3} in \cite{Ba1,Ba2} has its old vertices at depth $>D^\flat$, so \cite[Prop.~1]{Ba3} holds at $L=5$ in its printed form.
\end{proposition}
\begin{proof}
The side $r_\square$ enters \cite[pp.~433--434]{Ba3} in five places. (1) \emph{Containment:} \eqref{eq:a01:N} with $r_\square=r^\flat$. (2) \emph{Old-vertex classes lie in $\square_1$:} the expansion in $\tilde B'$ uses \cite[(3.14), (3.44), (3.45)]{Ba1} to order $\bar n$; only lines between localisations whose boxes meet stay internal \cite[p.~432]{Ba3}, so along an expanded line a new vertex advances at most $2$, and the new vertices lie within $2N_B$ of the class-component cube; the boxes of all graphs of a connected component lie in a grid cube of side $3\ell_{\bar n}$ concentric with one member's cube (``localizations \dots\ in agreement'', \cite[p.~432]{Ba3}); so the class lies in a cube of side $\ell^\sharp$, and $3\ell^\sharp\le r^\flat$. (3) \emph{The two $\delta G$ external fields are $O(1)$:} $\delta G_i(\Omega,\square,\tilde B)$ is evaluated at points of $\square_1$, at distance $\ge r^\flat\ge R_{\mathrm{Ba}}+12$ from $\square^c$, with sources in $\square_1$, so \cite[(2.26)]{Ba1} for the pair $(\square,\Omega)$ gives $\norm{h\,\delta G_i\,h'}_{1,\alpha}\le c_{\mathrm{Ba}}e^{-\delta_0r^\flat}$; for $\delta G_i(\square,\eta\Z^d,0)$ take a torus $T'\supset\square$ with side a multiple of $2s_b$ and at least $\max\{3s_b,3r^\flat+2s_b+2\}$: at zero background $G_i(T',0)$ is the periodisation of $G_i(\eta\Z^d,0)$, so $\delta G_i(\square,\eta\Z^d,0)=\delta G_i(\square,T',0)+\sum_{n\ne0}G_i(\eta\Z^d,0)(\cdot,\cdot+n\ell')$, the first term obeying \cite[(2.26)]{Ba1} on $T'$ with constants independent of $\ell'$ (the constants of \cite{Ba4} depend on $d$, $a$, $M$, $L$ only) and the second vanishing as $\ell'\to\infty$ by \cite[(2.25)]{Ba1}. Both are external fields of dimension $0$ \cite[(2.1), p.~422]{Ba3}, so an $O(1)$ bound is all that is used. (4) \emph{The constant field:} for $x\in\square$ and $c$ the centre, $|\tilde B(x)-\tilde B(c)|\le d(\frac32r^\flat+s_b)\sup_\square|\partial^\eta\tilde B|=O(r^\flat p_i)=O(p_i)$ with $B_0\defeq\tilde B(c)$, in place of the printed $O(r_ip_i)\le O(p_i^2)$; every use of this bound is an upper bound monotone in the size of $\tilde B'$ (vertices weighted by $e(\ell_i)|\tilde B'|$, the remainder \cite[(1.16)]{Ba3} by $(e(\ell_i)^{1-\alpha})^{n+n'}$, the norm comparison between the norms \cite[(1.32)]{Ba3} at $B_0$ and at $\tilde B$ costing $(1+O(p_i))^{\#\mathrm{legs}}$, absorbed by the slack $\ell_i^{(4-d)/4-\beta}$ of the applications \cite[p.~421]{Ba3}); so a smaller cube only helps. (5) \emph{Gauge-out of $B_0$:} $B_0=\partial^\eta\lambda_0$ with $\lambda_0(x)=B_0\cdot(x-c)$ on $\square$; the Neumann covariant Laplacian on the bonds of $\square$ and the averaging $Q_i$ (contours inside blocks contained in $\square$) transform covariantly, $G_i(\square,B_0)=U(\lambda_0)G_i(\square,0)U(\lambda_0)^*$, and every class is invariant when the external scalar fields are transformed simultaneously \cite[(2.69)--(2.74)]{Ba2}; exact as long as $\square$ does not wrap around the torus (no condition on the cube). (6) \emph{$\Omega$-freeness after the split:} every old vertex lies in $\square_1\subset\square\subset\Omega$ and every new vertex in $\square$, so the characteristic function of $\Omega$ is $1$ wherever a vertex or a line end occurs; every object depending on $\Omega$ or on the cube ($\delta G_i(\Omega,\Omega_2,\tilde B)$, \cite[(1.16)]{Ba3}, the two $\delta G$'s, the gauge factors) has been moved into external fields, used on \cite[pp.~434--445]{Ba3} only through the norms \cite[(3.3)]{Ba3} at points of $\square_1$; the $\Omega_1$-dependent objects $g_i$, $1_{\Omega_1}$, $\dm(e',g_i,\lambda',\Omega_1,x)$ stay in the vertices, fixed before the step and entering only identities between a graph and its counterterm with the same restriction; on \cite[pp.~434--445]{Ba3} every propagator is $G_i(0)$, a slice, $G^\xi$, $C^\xi$ or $\Pi$, every sum runs over $\eta\Z^d$, $\xi\Z^d$, the cubes $\Delta(v)$, $\Omega_1$ or the torus identified with $\eta\Z^d$, Taylor's formula \cite[(3.10)]{Ba3} is applied to legs and to $g,g'$, and the $\lambda$-vertices are matched against counterterms with the same localisation \cite[(3.19), pp.~444--445]{Ba3}: none of this refers to the cube side or to $\dist(\cdot,\Omega^c)$. (7) \emph{Chains without old vertices} \cite[p.~433]{Ba3}: a class with no old vertex is a path $v_1$--$\cdots$--$v_m$ of new vertices of one product $(G_0V)^mG_0$, $G_0\defeq G_i(\square,B_0)$, with the cut lines at its ends as external scalar fields; the printed route through \cite[Prop.~2.1]{Ba3} is not used (its summation by parts \cite[(2.8)]{Ba3} is printed on the torus and on the Neumann box produces a face term where the $\tilde B'$-localisation does not vanish); instead, with $\norm\psi_E^2\defeq\ip\psi{G_0^{-1}\psi}$, $G_0^{-1}\ge D_{B_0}^*D_{B_0}$ and $G_0^{-1}\ge\gamma_\square$ (constant field, gauged out), each vertex of \cite[(1.8)--(1.11), (1.13)--(1.15)]{Ba3} is a bounded form in $\norm\cdot_E$: the form of \cite[(1.8)]{Ba3} is $\kappa\sum_{b\subset\square}\eta^d(D_{B_0}\phi)(b)\cdot M(b)\chi(b_-)$ with $|M(b)|\le C(\sup|\tilde B'|)^{n'}$, bounded through $\norm{D_{B_0}\phi}\le\norm\phi_E$ and $\norm\chi\le\gamma_\square^{-1/2}\norm\chi_E$; \cite[(1.9), (1.11)]{Ba3} likewise, with $|R_{\bar n+1}|\le1$; \cite[(1.10)]{Ba3} carries no derivative and uses $\gamma_\square^{-1/2}$ on both legs; a block pair \cite[(1.13)--(1.14)]{Ba3} uses $|\tilde B'(\Gamma)|\le2d\sup|\tilde B'|$ on its unit-block contours and $\sum_{x\in B(y)}\eta^d|\phi(x)|\le\norm\phi$; every form norm is $\le C(e(\ell_i)\sup|\tilde B'|)^{n'_v}$. Hence $E(H)=\ip{G_0^{1/2}K_1^*\Phi_1}{\hat K_2\cdots\hat K_{m-1}G_0^{1/2}K_m\Phi_m}$, $\hat K_v\defeq G_0^{1/2}K_vG_0^{1/2}$, is bounded by the product of the form norms times the end norms $\norm{h\Phi}_E$, where $h=1$ on the $2\eta$-neighbourhood of the end cube, $|\partial^\eta h|\le c_h$, and $\norm{h\Phi}_E^2=\norm{D_{B_0}(h\Phi)}^2+m^2\ell_i^2\norm{h\Phi}^2+a_i\norm{Q_i(B_0)h\Phi}^2\le C(1+c_h)^2\norm{h\Phi}_{1,0}^2$ by the Leibniz rule on a cube of side $\le4$; so $\norm{1_{\square^\eta(v)}\Phi}+\norm{1_{\square^\eta(v)}D_{B_0}\Phi}\le C\norm{h\Phi}_{1,0}$ is all that is used of an end field. The cut lines are bounded pointwise on $\square$, with their first differences in each slot, at pairs of boxes at distance $\ge1-2\eta\ge\frac35$, by \cref{lem:a01:RBK}: $|D^a_{B_0,x}D^b_{B_0,x'}G_0(x,x')|\le Ce^{-\delta|x-x'|}$ for $|x-x'|\ge\frac35$ and $a,b\in\{0,1\}$, at every pair of points of $\square$, its faces included, with no depth condition; and a cut line $K$ contracted between two functionals $A$, $B$ bounded by the sup-and-difference norms of the two end regions obeys $|A(x\mapsto B(K(x,\cdot)))|\le\norm A\norm B\max_{a,b\in\{0,1\}}\sup|D^a_xD^b_{x'}K(x,x')|$, because $x\mapsto B(K(x,\cdot))$ is linear in $K$ and its differences are $B(D_xK(x,\cdot))$ bond by bond, so no hypothesis $\dist\ge1$ on the two supports is needed and distance exactly $1$ is covered. A cut line between a chain end and an old-vertex class is an external scalar field of that class, read through its norms \cite[(3.3)]{Ba3} after localisation to the $\eta$-neighbourhood of the class's cube by a function $h'$ as in \cref{lem:a01:RBK}(c), the chain slot anywhere in $\square$: by the same linearity, with $B$ the class functional, bounded by Ba{\l}aban's norm \cite[(1.32)]{Ba3} of its external field, $|A(x\mapsto B(h'K(x,\cdot)))|\le\norm A\norm B\max_{a\in\{0,1\}}\sup_x\norm{h'D^a_xK(x,\cdot)}$ with the norm \cite[(1.32)]{Ba3} in the last factor, and \cref{lem:a01:RBK}(c) bounds that factor by $C(1+c_h)^2$ at every $x$ of the chain-end region, the faces of $\square$ included. The vertex forms are sums over the bonds of $\square$, so vertices at the faces need no separate treatment. (8) \emph{The small factors of \cite{Ba2} survive:} the factors $O(e^{-\delta_1r_i})$ that \cite[p.~422]{Ba3} produces at domain changes come from $\exp[-\delta_0d(\{\square(v)\})]$ in \cite[(1.33)]{Ba3}, i.e.\ from the lines removed before the cube step, whose bound depends on $\Omega$ only; the two-domain inputs enter as external fields, linearly. So the only place where $r_\square=r_i$ was used is (1), where $r^\flat$ serves by hypothesis; the constant acquires the factors of (3) and (4), both of the printed kind. For the last claim: a vertex anchored at $\Omega_1$ has depth $\ge5r(\ell_{i-1})/5-3\ge r_i-3\ge R-3\ge D^\flat$ (\cref{lem:a01:X0}(a), $r(\ell_{i-1})\ge r_i\ge R$); a vertex anchored at $\supp\tilde A$ has depth $>2r_i-2-M_\theta-2\eta\ge2R-2.4-M_\theta>D^\flat$ (\cref{lem:a01:X0}(a)); the other steps of the proof need depth at most $R_{\mathrm{Ba}}+13<D^\flat$ at old vertices or within $2\sqrt3+1$ of $\supp\tilde A$.
\end{proof}
% a01_H0rel part 6x1, continued: two-slot bounds of the box propagator at a constant field (input from a01_p6x1.tex after prop:a01:P12)
\begin{lemma}[two-slot bounds of the box propagator at a constant field]\label{lem:a01:RBK}
Let $i\ge1$, $\eta\defeq L^{-i}$ the fine spacing in scale-$i$ units, $\square$ a box of large blocks (a product of intervals, each a union of large blocks) that does not wrap around the torus, with Neumann conditions on its faces, $B_0$ a constant field, and $G_0\defeq G_i(\square,B_0)$ the propagator of step (7) of the proof of \cref{prop:a01:P12}, $G_0^{-1}=D_{B_0}^*D_{B_0}+m^2\ell_i^2+a_iQ_i(B_0)^*Q_i(B_0)$ on $\ell^2(\square\cap\eta\Z^d)$ with the uniform block average $Q_i$. Kernels are taken for the measure $\eta^d$, and $D^a_{B_0,x}$, $a\in\{0,1\}$, denotes the identity or a covariant forward difference on a bond of $\square$ at $x$. \textup{(a)} There are $C,\delta>0$ depending on $(d,L,a)$ and the mass range only, free of $i$, $\eta$, $\square$, $B_0$ and the positions, such that for all $x,x'\in\square$ and $a,b\in\{0,1\}$ ($d\ge3$; here $d=3$)
\[
|D^a_{B_0,x}D^b_{B_0,x'}G_0(x,x')|\ \le\ C\,(\eta+|x-x'|)^{-(d-2+a+b)}\,e^{-\delta|x-x'|},
\qquad\text{so}\qquad |D^a_{B_0,x}D^b_{B_0,x'}G_0(x,x')|\le5C\,e^{-\delta|x-x'|}\ \text{ for }|x-x'|\ge\tfrac35\ (d=3).
\]
\textup{(b)} No depth condition enters: the bound holds at every pair of points of $\square$, the bonds of its faces included.

\textup{(c)} \emph{One slot in Ba{\l}aban's norm.} Let $\Delta,\Delta'$ be two cubes of $\square$ at distance $\ge1$, let $h'$ be a localisation function on $\square$, $0\le h'\le1$, equal to $1$ on the $\eta$-neighbourhood of $\Delta'$ and vanishing at distance $\ge\frac12$ from $\Delta'$, with $|\partial^\eta h'|\le c_h$ and $|\partial^\eta\partial^\eta h'|\le c_h$, and let $x$ lie within $\eta$ of $\Delta$. Then for $a\in\{0,1\}$ the field $\Psi\defeq h'(\cdot)\,D^a_{B_0,x}G_0(x,\cdot)$ satisfies the three clauses of Ba{\l}aban's external-field norm \cite[(1.32)]{Ba3} (sup, first differences, $\alpha_0$-H\"older quotients of the first differences on unit balls) on all of $\square$, its faces included, with the bound $C(1+c_h)^2$, $C$ of the same kind as above.
\end{lemma}
\begin{proof}
(1) \emph{Gauge-out.} With $\lambda_0(x)\defeq B_0\cdot(x-c)$ on the non-wrapping $\square$, $G_i(\square,B_0)=U(\lambda_0)G_i(\square,0)U(\lambda_0)^*$ and $D_{B_0}(U(\lambda_0)\phi)=U(\lambda_0)D\phi$ bond by bond (step (5) of the proof of \cref{prop:a01:P12}; $|U|=1$), so every modulus in the claim equals the one at $B_0=0$. Let $B_0=0$.

(2) \emph{Even images.} Let $\square=\prod_\mu[\alpha_\mu,\beta_\mu)$ and let $T_\square$ be the torus of periods $2(\beta_\mu-\alpha_\mu)$, each a multiple of $2s_b$, so that the large blocks of $\square$ tile $T_\square$ periodically; let $\mathcal W_\square\cong(\Z_2)^d$ be the group generated by the reflections in the face planes $x_\mu=\alpha_\mu-\eta/2$ of $\square$ (planes between sites on block boundaries; on $T_\square$ the reflection in $x_\mu=\beta_\mu-\eta/2$ is the same element). They map sites to sites, bonds to bonds and blocks to blocks, the uniform average commutes with them, and the mass term is constant; every bond of $T_\square$ that crosses a face plane joins a site to its mirror image, where a $\mathcal W_\square$-even field has zero difference, and the remaining bonds of $T_\square$ inside $\square$ are the bonds of $\square$. Hence, exactly as in \cref{lem:a01:PE}(a) and \cref{lem:a01:NL}, the Neumann operator on $\square$ is the even restriction of the periodic one, $G_i(T_\square,0)\Ebar=\Ebar\,G_i(\square,0)$ with $\Ebar$ the even extension, and
\[
G_i(\square,0)(x,x')=\sum_{g\in\mathcal W_\square}G_i(T_\square,0)(x,gx')\qquad(x,x'\in\square),
\]
a difference in $x'$ along a bond of $\square$ acting on the image terms along the reflected bonds. Every image satisfies $\dist_{T_\square}(x,gx')\ge|x-x'|$: coordinatewise, for $x_\mu,x'_\mu\in[\alpha_\mu,\beta_\mu-\eta]$ and $r$ the reflection in $\beta_\mu-\eta/2$, $rx'_\mu-x_\mu=|x_\mu-x'_\mu|+2(\beta_\mu-\tfrac\eta2-\max(x_\mu,x'_\mu))\ge|x_\mu-x'_\mu|$, and the complementary distance around the period is $x_\mu+x'_\mu-2\alpha_\mu+\eta\ge|x_\mu-x'_\mu|$.

(3) \emph{Periodisation.} At zero field $G_i(T_\square,0)(x,y)=\sum_{n\in\Z^d}G_i(\eta\Z^d,0)(x,y+2\ell n)$, $\ell\defeq(\beta_\mu-\alpha_\mu)_\mu$ (the periodic operator is the periodisation of the free-lattice one, as in step (3) of the proof of \cref{prop:a01:P12}), every term having $|x-y-2\ell n|\ge\dist_{T_\square}(x,y)$, and at most $(2+k)^d$ values of $n$ have $|x-y-2\ell n|<k+1$, every period being $\ge1$. So, with the free-lattice bound (4) below at the rate $\delta'$ and $\rho\defeq|x-x'|$,
\[
|D^a_xD^b_{x'}G_i(\square,0)(x,x')|\ \le\ 2^dC'(\eta+\rho)^{-(d-2+a+b)}\sum_{k\ge\lfloor\rho\rfloor}(2+k)^de^{-\delta'k}\ \le\ C(\eta+\rho)^{-(d-2+a+b)}e^{-\delta'\rho/2}.
\]

(4) \emph{The free lattice.} On $\eta\Z^d$ at zero field, $G_i(\eta\Z^d,0)=G_1+\sum_{j=1}^{i-1}G_{(j)}$ with the slices $G_{(j)}\defeq G_{j+1}-G_j=a_jG_jQ^*C^{(j)}QG_j$ (the Schur identity behind \cref{lem:a01:legs}(b) and \cite[(2.17)]{K1}, with $C^{(j)}$ as in \cref{lem:a01:A23}(c)), so that the sum telescopes to $G_i$. In unit-$j$ units ($j$-blocks of side one, spacing $L^{-j}$) let $\ell_j(\cdot,u)\defeq a_jG_jQ^*e_u$ be the leg of the unit datum on the $j$-block $u$; since $G_j$ is real symmetric at zero field, $G_{(j)}(y,y')=a_j^{-1}\sum_{u,v}\ell_j(y,u)C^{(j)}(u,v)\ell_j(y',v)$, each variable sitting in its own leg, so
\[
D^a_yD^b_{y'}G_{(j)}(y,y')=a_j^{-1}\sum_{u,v}(D^a\ell_j)(y,u)\,C^{(j)}(u,v)\,(D^b\ell_j)(y',v).
\]
The legs obey $|D^a\ell_j(y,u)|\le Ce^{-\delta_2\dist(y,B(u))}$ at every site and every bond of the free lattice, $a\in\{0,1\}$ (the leg bounds \cite[(3.64)]{K1}, \cite[(2.12)]{Ba3}, at the uniform average $f=1$ in \cref{prop:a02:FL}; on the free lattice every bond is interior, so they are also \cref{lem:a01:A4}(b) and the interior case (4) of the proof of \cref{lem:a01:A5}), and $|C^{(j)}(u,v)|\le Ce^{-\delta_1|u-v|}$ (\cite[Thm.~3.3, (3.6)]{K1}; \cref{lem:a01:A23}(c)). With $|y-y'|\le\dist(y,B(u))+|u-v|+\dist(y',B(v))+\sqrt d$ and two geometric sums over $u,v\in\Z^d$, $|D^a_yD^b_{y'}G_{(j)}(y,y')|\le C_Se^{-\delta_S|y-y'|}$ at every pair of points, $\delta_S\defeq\frac12\min\{\delta_1,\delta_2\}$. In scale-$i$ units the slice $j$ has the kernel $s_j^{2-d}G_{(j)}(x/s_j,x'/s_j)$ with $s_j\defeq L^{j-i}$, and a scale-$i$ difference is $s_j^{-1}$ times a unit-$j$ one, so slice $j$ contributes at most $C_Ss_j^{-(d-2+a+b)}e^{-\delta_S|x-x'|/s_j}$; $G_1$, read in unit-$1$ units at the fixed spacing $L^{-1}$, obeys the same bound with $s\defeq L\eta$ and a constant depending on $L$ (its kernel is $L^d\ip{e_y}{G_1e_{y'}}$, bounded with the decay of \cref{lem:a01:A23}(a), and a difference at spacing $L^{-1}$ costs a factor $2L$). Finally, for $p\defeq d-2+a+b>0$ and $s$ ranging over the powers of $L$ in $[L\eta,L^{-1}]$: $e^{-\delta_S\rho/s}\le e^{-\delta_S\rho/2}e^{-\delta_S\rho/(2s)}$ as $s\le1$; the terms with $s\ge\max(\rho,\eta)$ sum to at most $(1-L^{-p})^{-1}\max(\rho,\eta)^{-p}$, and the terms with $s<\rho$ are $\rho^{-p}(\rho/s)^pe^{-\delta_S\rho/(2s)}$ with $\rho/s$ running over powers of $L$, summable to at most $C_p\rho^{-p}$. Hence $|D^a_xD^b_{x'}G_i(\eta\Z^d,0)(x,x')|\le C'(\eta+\rho)^{-p}e^{-\delta'\rho}$ with $\delta'\defeq\delta_S/2$, which is (4), and (3) gives the claim with $\delta\defeq\delta'/2$. The last clause of the statement is $(\eta+\rho)^{-p}\le(5/3)^3<5$ for $\rho\ge\frac35$.

(5) \emph{The clause (c).} By (1) every modulus is the one at $B_0=0$. The reflections of (2) and the translations of (3) are isometries of the lattice mapping bonds to bonds, so each clause of $x'\mapsto D^a_xG_i(\square,0)(x,x')$ at a pair $x',y'\in\square$ is bounded by the sum over the images of the same clause of $x'\mapsto D^a_xG_i(\eta\Z^d,0)(x,x')$ at the pair $(gx'+2\ell n,gy'+2\ell n)$, whose points are at least as far from $x$ as $x'$, $y'$ are, with the count of (3). Write $\Psi_0(x')\defeq D^a_xG_i(\eta\Z^d,0)(x,x')$ and $\rho'\defeq\min\{|x'-x|,|y'-x|\}$. \emph{Claim:} on the free lattice at zero field the sup and first-difference clauses of $\Psi_0$ at $x'$, and the H\"older quotient at a pair $x',y'$ with $|x'-y'|\le1$, are $\le Ce^{-\delta\rho'}$ whenever $\rho'\ge\frac1{20}$. Indeed, by the factorisation of (4), in unit-$j$ units $D^a_xG_{(j)}(x,\cdot)=\sum_v\ell_j(\cdot,v)\,a_j^{-1}\beta_j(v)$ is the leg of the block data $a_j^{-1}\beta_j$, $\beta_j(v)\defeq\sum_u(D^a\ell_j)(x,u)C^{(j)}(u,v)$, with $|\beta_j(v)|\le Ce^{-\delta_S\dist(x,B(v))}$ by the leg and $C^{(j)}$ bounds of (4) and one geometric sum. Split the data into the blocks with $\dist(x,B(v))\le\rho'_j/2$, $\rho'_j\defeq\rho'/s_j$ the distance in unit-$j$ units, and the rest. The leg of bounded block data obeys the sup bound of \cref{lem:a01:A4}(b), the first-difference bound of \cref{lem:a01:A5} and its H\"older bound on balls of radius $\frac1{16}$ with decay $e^{-\delta_3\dist(\cdot,S)}$ from the support $S$ of the data (the interior case, every bond of the free lattice being interior; printed: \cite[(3.63)--(3.65)]{K1}, \cite[(2.12)]{Ba3}); the far part has data of size $\le Ce^{-\delta_S\rho'_j/2}$, and the near part has its support within $\rho'_j/2$ of $x$, so at distance $\ge\rho'_j/2$ from $x'$ and $y'$; a pair at unit-$j$ distance $>\frac1{16}$ is bounded by $2\cdot16^{\alpha_0}$ times the first differences. So in unit-$j$ units every clause of slice $j$ is $\le Ce^{-\delta_6\rho'_j}$, $\delta_6\defeq\frac14\min\{\delta_S,\delta_3\}$. In scale-$i$ units the sup clause of slice $j$ scales as $s_j^{-(d-2+a)}$, the first differences as $s_j^{-(d-1+a)}$ and the H\"older quotients as $s_j^{-(d-1+a+\alpha_0)}$ (a scale-$i$ pair at distance $t\le1$ is a unit-$j$ pair at distance $t/s_j$; for $t>s_j/16$ use two first differences, $2\cdot16^{\alpha_0}s_j^{-(d-1+a)}t^{-\alpha_0}\le2\cdot16^{\alpha_0}s_j^{-(d-1+a+\alpha_0)}$); the slice $G_1$ at the fixed spacing contributes a bounded multiple of its first differences to the H\"older clause, every pair being at distance $\ge\eta$. The sum over $s$ is the geometric sum of (4) with $d-1+a+\alpha_0$ in place of $p$, giving $\le C(\eta+\rho')^{-(d-1+a+\alpha_0)}e^{-\delta'\rho'}\le Ce^{-\delta'\rho'}$ for $\rho'\ge\frac1{20}$; this is the claim, and with the image count it holds for $x'\mapsto D^a_xG_i(\square,0)(x,x')$ at pairs of $\square$ with $\rho'\ge\frac1{20}$. \emph{Localisation.} $\Psi=h'\Psi_\square$, $\Psi_\square\defeq D^a_xG_i(\square,0)(x,\cdot)$, and $D(h'\Psi_\square)=h'(\cdot+\eta e)D\Psi_\square+(\partial^\eta h')\Psi_\square$ bond by bond. Both $\Psi$ and $D\Psi$ vanish outside $S'\defeq\{\dist(\cdot,\Delta')\le\frac12+\eta\}$, whose points are at distance $\ge\frac12-2\eta\ge\frac1{10}$ from $x$ ($\eta\le\frac15$, $x$ within $\eta$ of $\Delta$, $\dist(\Delta,\Delta')\ge1$). Hence the sup and first-difference clauses of $\Psi$ are those of $\Psi_\square$ on $S'$ times $1+c_h$. For a H\"older pair $x',y'$ with both points at distance $\ge\frac1{20}$ from $x$, the Leibniz rule and the bounds $|h'(x')-h'(y')|\le dc_h|x'-y'|$, $|\partial^\eta h'(x')-\partial^\eta h'(y')|\le dc_h|x'-y'|$ give the quotient $\le C(1+c_h)^2$ from the three clauses of $\Psi_\square$; a pair with one point $y'$ within $\frac1{20}$ of $x$ has $D\Psi(y')=0$ and $|x'-y'|\ge\frac1{10}-\frac1{20}=\frac1{20}$ whenever $D\Psi(x')\ne0$, so its quotient is $\le20^{\alpha_0}$ times a first difference. This is (c).
\end{proof}

On the cube the same device, a constant box $Q_x$ of large blocks around each vertex cube at depth $\ge R_*$, is step (1) of the proof of \cref{thm:a01:B2}; the depth supply is \cref{lem:a01:X0}(a) there too.
